\PassOptionsToPackage{table}{xcolor}
\documentclass[10pt,aspectratio=169]{beamer}
\newcommand{\LectureNumber}{23}
\input{course-theme.tex}
\title{Ragged and multidimensional arrays}
\subtitle{CSC103 Programming Fundamentals / Lecture 23}
\author{Dr. Muhammad Farhan}
\institute{Associate Professor of Computer Science\\Department of Computer Science\\COMSATS University Islamabad\\Sahiwal Campus}
\date{Fall 2026}
\begin{document}
\begin{frame}\titlepage\end{frame}

\begin{frame}[fragile,t]{The problem for this lecture}
\begin{columns}[T,onlytextwidth]\begin{column}{0.60\textwidth}
Write a method to count all elements in a non-null ragged array with non-null rows. Test \{\{1\},\{2,3\},\{\}\}.
\end{column}\begin{column}{0.35\textwidth}\small
Assume the outer array and every row are non-null.
\end{column}\end{columns}
\note{Introduce the target problem before explaining the tools. Revisit it in the guided solution later.\par Syllabus reading: Deitel Ch 7. CLO-2.}
\end{frame}

\begin{frame}[fragile,t]{Learning objectives}
\begin{columns}[T,onlytextwidth]\begin{column}{0.60\textwidth}
\begin{itemize}
\item Pass a two-dimensional array to a method
\item Handle ragged rows
\item Explain copying and additional dimensions
\end{itemize}
\end{column}\begin{column}{0.35\textwidth}\small
CLO-2

Two-dimensional parameters\par Ragged arrays\par Shallow and deeper copying\par Additional dimensions
\end{column}\end{columns}
\note{Explain how each objective supports the opening problem. The final questions check these objectives.\par Syllabus reading: Deitel Ch 7. CLO-2.}
\end{frame}

\begin{frame}[fragile,t]{Two-dimensional parameters}
\begin{itemize}
\item A method can accept int[][] in its parameter list.
\item As with other arrays, the reference value is copied.
\item Element changes can affect caller-visible data.
\end{itemize}
\vspace{1mm}\begin{block}{Example}
\begin{lstlisting}
static int total(int[][] data) {
  int sum = 0;
  for (int[] row : data)
    for (int v : row) sum += v;
  return sum;
}
\end{lstlisting}
\end{block}
\note{State a contract: data and each row are non-null. Enhanced loops work for empty and unequal-length rows.\par Syllabus reading: Deitel Ch 7. CLO-2.}
\end{frame}

\begin{frame}[fragile,t]{Ragged arrays}
\begin{itemize}
\item Rows may have different lengths.
\item Allocate the outer array and each row separately when needed.
\item A newly created outer reference array starts with null rows.
\end{itemize}
\vspace{1mm}\begin{block}{Example}
\begin{lstlisting}
int[][] groups = new int[2][];
groups[0] = new int[]{1};
groups[1] = new int[]{2, 3, 4};
\end{lstlisting}
\end{block}
\note{groups[0].length=1 and groups[1].length=3. Accessing a row before assigning it would dereference null.\par Syllabus reading: Deitel Ch 7. CLO-2.}
\end{frame}

\begin{frame}[fragile,t]{Shallow and deeper copying}
\begin{itemize}
\item Copying only the outer array preserves shared row references.
\item Changing a shared row element affects both structures.
\item Independent rows require copying each row too.
\end{itemize}
\vspace{1mm}\begin{block}{Example}
\begin{lstlisting}
int[][] copy = original.clone();
copy[0][0] = 99;
// original[0][0] is now 99
\end{lstlisting}
\end{block}
\note{For primitive two-dimensional arrays, clone each row to copy all element storage. Deeper object structures can require additional decisions.\par Syllabus reading: Deitel Ch 7. CLO-2.}
\end{frame}

\begin{frame}[fragile,t]{Additional dimensions}
\begin{itemize}
\item Each extra dimension introduces another array level.
\item Give every index a domain meaning.
\item Prefer the smallest number of dimensions that expresses the task.
\end{itemize}
\vspace{1mm}\begin{block}{Example}
\begin{lstlisting}
int[][][] attendance = new int[2][3][4];
// [section][week][session]
\end{lstlisting}
\end{block}
\note{There are 24 int elements in this rectangular allocation. A null or ragged subarray is possible in other constructions, so do not assume a flat contiguous block.\par Syllabus reading: Deitel Ch 7. CLO-2.}
\end{frame}

\begin{frame}[fragile,t]{A ragged array needs a row-specific bound}
\begin{itemize}
\item The outer length counts row references, not every stored element.
\item Each non-null row supplies its own length.
\item Adding the row lengths counts all primitive elements, including zero for an empty row.
\end{itemize}
\vspace{1mm}\begin{block}{Example}
\begin{lstlisting}
{{1},{2,3},{}} has row lengths 1,2,0.
Element count = 1 + 2 + 0 = 3.
\end{lstlisting}
\end{block}
\note{Explain the mechanism using the displayed example. The outer length counts row references, not every stored element. Each non-null row supplies its own length. Adding the row lengths counts all primitive elements, including zero for an empty row.\par Syllabus reading: Deitel Ch 7. CLO-2.}
\end{frame}

\begin{frame}[fragile,t]{Copy depth follows the structure}
\begin{itemize}
\item Cloning a two-dimensional array copies the outer row references.
\item Cloning each row as well produces separate primitive element storage.
\item A deeper structure may require further copying decisions.
\end{itemize}
\vspace{1mm}\begin{block}{Example}
\begin{lstlisting}
int[][] copy = new int[a.length][];
for (int r=0; r<a.length; r++)
  copy[r] = a[r].clone();
\end{lstlisting}
\end{block}
\note{Explain the mechanism using the displayed example. Cloning a two-dimensional array copies the outer row references. Cloning each row as well produces separate primitive element storage. A deeper structure may require further copying decisions.\par Syllabus reading: Deitel Ch 7. CLO-2.}
\end{frame}


\begin{frame}[fragile,t]{Ragged rows have different lengths}
\begin{center}\begin{tikzpicture}
\node[text=accent,font=\small] at (-1.7,0.0) {row 0};
\node[box,text width=13mm] (r0c0) at (0.0,0.0) {8};
\node[text=accent,font=\small] at (-1.7,-1.2) {row 1};
\node[box,text width=13mm] (r1c0) at (0.0,-1.2) {2};
\node[box,text width=13mm] (r1c1) at (1.8,-1.2) {4};
\node[box,text width=13mm] (r1c2) at (3.6,-1.2) {6};
\node[text=accent,font=\small] at (-1.7,-2.4) {row 2};
\node[hot] (empty) at (1.8,-2.4) {Empty row};
\end{tikzpicture}\end{center}
\begin{block}{What to notice}The outer length counts rows; each row supplies its own element count.\end{block}
\note{Trace each labelled connection. The outer length counts rows; each row supplies its own element count.}
\end{frame}
\begin{frame}[fragile,t]{Key distinctions}
\small\renewcommand{\arraystretch}{1.5}
\rowcolors{2}{pale}{white}
\begin{tabularx}{\textwidth}{>{\raggedright\arraybackslash}X>{\raggedright\arraybackslash}X>{\raggedright\arraybackslash}X>{\raggedright\arraybackslash}X}
\rowcolor{navy}
\color{white}\bfseries Row & \color{white}\bfseries Contents & \color{white}\bfseries Length & \color{white}\bfseries Valid indexes\\
0 & \{1\} & 1 & 0\\
1 & \{2,3\} & 2 & 0,1\\
2 & \{\} & 0 & None\\
\end{tabularx}
\vspace{3mm}\footnotesize These distinctions determine which operation or control structure fits the problem.
\note{\par Syllabus reading: Deitel Ch 7. CLO-2.}
\end{frame}

\begin{frame}[fragile,t]{First example: methods}
\begin{lstlisting}
static int total(int[][] data) {
  int sum = 0;
  for (int[] row : data)
    for (int v : row) sum += v;
  return sum;
}
\end{lstlisting}
\note{Method definitions belong inside the class and outside main. Explain the parameters and return values.\par Syllabus reading: Deitel Ch 7. CLO-2.}
\end{frame}

\begin{frame}[fragile,t]{First example: execution}
\begin{lstlisting}
int[][] data = {{1, 2}, {3}};
System.out.println(total(data));
int[][] copy = data.clone();
copy[0][0] = 9;
System.out.println(data[0][0]);
\end{lstlisting}
\note{This is the main body from Java\_Examples/Lecture23.java. The source file includes the complete class. Predict the result before the next slide.\par Syllabus reading: Deitel Ch 7. CLO-2.}
\end{frame}

\begin{frame}[fragile,t]{First example: result}
\begin{columns}[T,onlytextwidth]\begin{column}{0.60\textwidth}
6\par 9\par The total method processes unequal row lengths.\par clone copied the outer references, so the first row remains shared.
\end{column}\begin{column}{0.35\textwidth}\small
Shallow and deeper copying

Additional dimensions
\end{column}\end{columns}
\note{Expected output and interpretation: 6
9
The total method processes unequal row lengths.
clone copied the outer references, so the first row remains shared.\par Syllabus reading: Deitel Ch 7. CLO-2.}
\end{frame}

\begin{frame}[fragile,t]{Guided practice}
\begin{columns}[T,onlytextwidth]\begin{column}{0.60\textwidth}
Write a method to count all elements in a non-null ragged array with non-null rows. Test \{\{1\},\{2,3\},\{\}\}.
\end{column}\begin{column}{0.35\textwidth}\small
Attempt the solution before continuing.

Use the definitions and examples from this lecture.
\end{column}\end{columns}
\note{Give students 8 minutes to attempt the problem. The following slides provide a complete solution. Original solution guidance: Start count=0 and add row.length for each row. The result is 3. The empty row contributes zero.\par Syllabus reading: Deitel Ch 7. CLO-2.}
\end{frame}

\begin{frame}[fragile,t]{Solution strategy}
\begin{columns}[T,onlytextwidth]\begin{column}{0.60\textwidth}
\begin{itemize}
\item Assume the outer array and every row are non-null.
\item Visit each row and add its length to a counter.
\item Return the accumulated count.
\end{itemize}
\end{column}\begin{column}{0.35\textwidth}\small
The next program uses fixed inputs so its result can be reproduced.
\end{column}\end{columns}
\note{Discuss why each step is needed. State the input assumptions before executing the program.\par Syllabus reading: Deitel Ch 7. CLO-2.}
\end{frame}

\begin{frame}[fragile,t]{Complete solution}
\begin{lstlisting}
public class Solution23 {
  public static void main(String[] args) throws Exception {
    int[][] values = {{1}, {2, 3}, {}};
    System.out.println(elementCount(values));
  }
  static int elementCount(int[][] values) {
    int count = 0;
    for (int[] row : values) {
      count += row.length;
    }
    return count;
  }
}
\end{lstlisting}
\note{Complete runnable file: Java\_Examples/Solution23.java. Inputs are fixed in the program.  \par Syllabus reading: Deitel Ch 7. CLO-2.}
\end{frame}

\begin{frame}[fragile,t]{Execution trace}
\small\renewcommand{\arraystretch}{1.5}
\rowcolors{2}{pale}{white}
\begin{tabularx}{\textwidth}{>{\raggedright\arraybackslash}X>{\raggedright\arraybackslash}X>{\raggedright\arraybackslash}X>{\raggedright\arraybackslash}X}
\rowcolor{navy}
\color{white}\bfseries Row visited & \color{white}\bfseries Length added & \color{white}\bfseries Count before & \color{white}\bfseries Count after\\
0 & 1 & 0 & 1\\
1 & 2 & 1 & 3\\
2 & 0 & 3 & 3\\
\end{tabularx}
\vspace{3mm}\footnotesize Follow the changing state in execution order.
\note{Manually derived trace for the complete solution. Assume the outer array and every row are non-null. Visit each row and add its length to a counter. Return the accumulated count.\par Syllabus reading: Deitel Ch 7. CLO-2.}
\end{frame}

\begin{frame}[fragile,t]{Solution result}
\begin{columns}[T,onlytextwidth]\begin{column}{0.60\textwidth}
3
\end{column}\begin{column}{0.35\textwidth}\small
Why this result follows

Return the accumulated count.
\end{column}\end{columns}
\note{Expected result: 3\par Syllabus reading: Deitel Ch 7. CLO-2.}
\end{frame}

\begin{frame}[fragile,t]{A common incorrect approach}
int[][] a = new int[3][];\par a[0][0] = 5;
\vspace{1mm}\begin{block}{Example}
\begin{lstlisting}
a[0] is initially null. Allocate the row, for example a[0]=new int[2], before assigning an element.
\end{lstlisting}
\end{block}
\note{Ask students to predict the failure before discussing the explanation. The right side gives the correction.\par Syllabus reading: Deitel Ch 7. CLO-2.}
\end{frame}

\begin{frame}[fragile,t]{Boundary and failure checks}
\begin{columns}[T,onlytextwidth]\begin{column}{0.60\textwidth}
Check the stated input assumptions.

Write a method that makes an independent copy of a ragged int[][] with non-null rows. Demonstrate that mutation of the copy leaves the original unchanged.
\end{column}\begin{column}{0.35\textwidth}\small
Start count=0 and add row.length for each row. The result is 3. The empty row contributes zero.
\end{column}\end{columns}
\note{Use the specification to derive each expected result. Exercise solution: Start count=0 and add row.length for each row. The result is 3. The empty row contributes zero.\par Syllabus reading: Deitel Ch 7. CLO-2.}
\end{frame}

\begin{frame}[fragile,t]{Check your understanding}
\begin{columns}[T,onlytextwidth]\begin{column}{0.60\textwidth}
Does outer-array clone duplicate each row? How many values are in new int[2][3][4]?
\end{column}\begin{column}{0.35\textwidth}\small
Write a prediction and explain the rule that supports it.
\end{column}\end{columns}
\note{Answer: No. It shares the row references. The rectangular allocation contains 24 int values.\par Syllabus reading: Deitel Ch 7. CLO-2.}
\end{frame}

\begin{frame}[fragile,t]{Answer and explanation}
\begin{columns}[T,onlytextwidth]\begin{column}{0.60\textwidth}
No. It shares the row references. The rectangular allocation contains 24 int values.
\end{column}\begin{column}{0.35\textwidth}\small
a[0] is initially null. Allocate the row, for example a[0]=new int[2], before assigning an element.
\end{column}\end{columns}
\note{Reconnect the answer to the relevant definition rather than treating it as a fact to memorise.\par Syllabus reading: Deitel Ch 7. CLO-2.}
\end{frame}


\begin{frame}[fragile,t]{Compare cases and predict outcomes}
\small\renewcommand{\arraystretch}{1.5}\rowcolors{2}{pale}{white}
\begin{tabularx}{\textwidth}{>{\raggedright\arraybackslash}X>{\raggedright\arraybackslash}X>{\raggedright\arraybackslash}X}\rowcolor{navy}
\color{white}\bfseries Row & \color{white}\bfseries Values & \color{white}\bfseries Length\\
0 & 8 & 1\\
1 & 2, 4, 6 & 3\\
2 & No elements & 0\\
\end{tabularx}\par\vspace{3mm}\begin{exampleblock}{Discuss}Explain each expected result before running the example. Which case would reveal a wrong assumption?\end{exampleblock}
\note{Read the first two columns and invite predictions before discussing the outcome.}
\end{frame}

\begin{frame}[fragile,t]{Apply the idea}
\begin{alertblock}{Independent practice}For \{\{8\},\{2,4,6\},\{\}\}, compute the total number of elements and their sum. Assume all rows are non-null.\end{alertblock}\vspace{3mm}\begin{enumerate}\item Work through the input by hand.\item Write or correct the program.\item Justify the expected result.\end{enumerate}\note{Allow 3–5 minutes, then compare answers in pairs. For \{\{8\},\{2,4,6\},\{\}\}, compute the total number of elements and their sum. Assume all rows are non-null.}
\end{frame}

\begin{frame}[fragile,t]{Solution and reasoning}
\begin{lstlisting}
int[][] data = {{8}, {2, 4, 6}, {}};
int count = 0, sum = 0;
for (int[] row : data) {
  count += row.length;
  for (int value : row) sum += value;
}
System.out.println(count + " " + sum);
\end{lstlisting}\vspace{1mm}\small\begin{itemize}
\item The lengths add to 4.
\item The values add to 20.
\item The empty row contributes zero to both totals without needing a special loop branch.
\end{itemize}\note{Answer: The lengths add to 4. The values add to 20. The empty row contributes zero to both totals without needing a special loop branch.}
\end{frame}

\begin{frame}[fragile,t]{Selecting the row with the largest sum}
\begin{itemize}
\item Compute one row sum at a time using that row's length.
\item Initialise the best sum from an actual row rather than zero if negatives are allowed.
\item A strict \textgreater{} update keeps the first row when sums tie.
\end{itemize}
\vspace{3mm}\footnotesize\textcolor{accent}{Source: supplied ISB campus slides. See speaker notes and ISB\_Source\_Map.md.}
\note{Adapted from the supplied ISB campus folder: Multi{-}dimensional arrays.pptx, slides 10, 66, 69. Adapted explanation and runnable implementation; sample inputs and traces added for this course.}
\end{frame}

\begin{frame}[fragile,t]{Worked problem: Selecting the row with the largest sum}
\begin{block}{Task}Find the row with the largest sum in the source matrix \{\{10,20,30\},\{40,50,60\},\{70,80,90\}\}.\end{block}
\small Input: Values are fixed in the program for a reproducible demonstration.\par\vspace{3mm}\begin{exampleblock}{Before running}Predict the output and identify the variables that change.\end{exampleblock}
\note{Adapted from the supplied ISB campus folder: Multi{-}dimensional arrays.pptx, slides 10, 66, 69. Adapted explanation and runnable implementation; sample inputs and traces added for this course.}
\end{frame}

\begin{frame}[fragile,t]{Complete ISB example (1/1)}
\begin{lstlisting}
public class IsbLecture23 {
  public static void main(String[] args) throws Exception {
    int[][] data = {{10,20,30}, {40,50,60}, {70,80,90}};
    int bestRow = 0;
    long bestSum = Long.MIN_VALUE;
    for (int row = 0; row < data.length; row++) {
      long sum = 0;
      for (int value : data[row]) sum += value;
      if (sum > bestSum) { bestSum = sum; bestRow = row; }
    }
    System.out.println("Row " + bestRow + ", sum " + bestSum);
  }

}
\end{lstlisting}
\note{Adapted from the supplied ISB campus folder: Multi{-}dimensional arrays.pptx, slides 10, 66, 69. Adapted explanation and runnable implementation; sample inputs and traces added for this course. Complete file: ISB\_Java\_Examples/IsbLecture23.java.}
\end{frame}

\begin{frame}[fragile,t]{Worked trace and expected result}
\small\renewcommand{\arraystretch}{1.4}\rowcolors{2}{pale}{white}
\begin{tabularx}{\textwidth}{>{\raggedright\arraybackslash}X>{\raggedright\arraybackslash}X>{\raggedright\arraybackslash}X}\rowcolor{navy}
\color{white}\bfseries Row index & \color{white}\bfseries Row sum & \color{white}\bfseries Best row afterward\\
0 & 60 & 0\\
1 & 150 & 1\\
2 & 240 & 2\\
\end{tabularx}\par\vspace{2mm}
\begin{block}{Expected console output}\ttfamily\footnotesize Row 2, sum 240\end{block}
\note{Adapted from the supplied ISB campus folder: Multi{-}dimensional arrays.pptx, slides 10, 66, 69. Adapted explanation and runnable implementation; sample inputs and traces added for this course. Trace and expected values independently worked for this edition.}
\end{frame}

\begin{frame}[fragile,t]{Extend the ISB example}
\begin{alertblock}{Practice}How would you extend the contract to an empty outer array and null rows?\end{alertblock}\vspace{3mm}\begin{exampleblock}{Explain your reasoning}Reject an empty outer array because there is no winning row. Reject null rows or define a separate policy. The demonstrated input is nonempty with non{-}null rows.\end{exampleblock}
\note{Adapted from the supplied ISB campus folder: Multi{-}dimensional arrays.pptx, slides 10, 66, 69. Adapted explanation and runnable implementation; sample inputs and traces added for this course. Answer: Reject an empty outer array because there is no winning row. Reject null rows or define a separate policy. The demonstrated input is nonempty with non{-}null rows.}
\end{frame}
\begin{frame}[fragile,t]{Lecture summary}
\begin{columns}[T,onlytextwidth]\begin{column}{0.60\textwidth}
\begin{itemize}
\item a two-dimensional array to a method
\item ragged rows
\item copying and additional dimensions
\end{itemize}
\end{column}\begin{column}{0.35\textwidth}\small
\begin{itemize}
\item Assume the outer array and every row are non-null.
\item Visit each row and add its length to a counter.
\item Return the accumulated count.
\end{itemize}
\end{column}\end{columns}
\note{Ask students to give one concrete example for the concepts listed. Use the worked solution to connect the topics.\par Syllabus reading: Deitel Ch 7. CLO-2.}
\end{frame}

\begin{frame}[fragile,t]{After-class practice}
\begin{columns}[T,onlytextwidth]\begin{column}{0.60\textwidth}
Write a method that makes an independent copy of a ragged int[][] with non-null rows. Demonstrate that mutation of the copy leaves the original unchanged.
\end{column}\begin{column}{0.35\textwidth}\small
Syllabus reading\par Deitel Ch 7

Next lecture: Exception handling
\end{column}\end{columns}
\note{Reading labels follow the supplied outline. Check topic headings against the available textbook edition. Require a program and expected results for the practice task.\par Syllabus reading: Deitel Ch 7. CLO-2.}
\end{frame}
\end{document}
